% part: intuitionistic-logic
% chapter: soundness-completeness
% section: completeness-thm

\documentclass[../../../include/open-logic-section]{subfiles}

\begin{document}

\olfileid{int}{sc}{cpl}

\olsection{પૂર્ણતા પ્રમેય}

\begin{thm}\ollabel{thm:completeness}
  જો $\Gamma \Entails !A$ હોય, તો $\Gamma \Proves !A$.
\end{thm}

\begin{proof}
  પ્રતિપક્ષી વિધાન સાબિત કરીએ: ધારો કે $\Gamma \Proves/ !A$.
  ત્યારે \olref[lin]{lem:lindenbaum} મુજબ એવો અભાજ્ય ગણ
  $\Gamma^* \supseteq \Gamma$ છે કે $\Gamma^* \Proves/ !A$.
  \olref[mod]{defn:canonical-model}માં વ્યાખ્યાયિત
  $\Gamma^*$ માટેનું કૅનોનિકલ નિદર્શ~$\mModel{M(\Gamma^*)}$
  લો. દરેક $!B \in \Gamma$ માટે $\Gamma^* \Proves !B$.
  નોંધો કે $\Gamma^*(\emptyseq) = \Gamma^*$.
  સત્યતા સહાયક પ્રમેય (\olref[tru]{lem:truth}) મુજબ
  બધા $!B \in \Gamma$ માટે
  $\mSat{M(\Gamma^*)}{!B}[\emptyseq]$ અને
  $\mSat/{M(\Gamma^*)}{!A}[\emptyseq]$ મળે છે.
  આથી $\Gamma \Entails/ !A$.
\end{proof}

\begin{prob}
  જો $!A$માં ફક્ત !!{propositional variable}s, $\lor$
  અને~$\land$ હોય, તો $\Entails/ !A$ છે તે બતાવો.
  આ પરથી સ્ફુરણાવાદી તર્કમાં $\lor$ અને~$\land$
  વડે $\lif$ વ્યાખ્યાયિત થઈ શકતું નથી તે તારવો.
\end{prob}

\begin{prob}
  પૂર્ણતા પ્રમેય વાપરીને સાબિત કરો કે જો
  $\Proves !A \lor !B$ હોય, તો $\Proves !A$ અથવા
  $\Proves !B$. (સંકેત: $\mSat/{M_1}{!A}$ અને
  $\mSat/{M_2}{!B}$ ધારો અને એવું નવું નિદર્શ
  $\mModel{M}$ રચો કે $\mSat/{M}{!A \lor !B}$.)
\end{prob}

\begin{prob}
  જો $\mModel{M}$ રેખીય ક્રમ વાપરતું સંબંધાત્મક
  નિદર્શ હોય, તો
  $\mSat{M}{(!A \lif !B)\lor(!B \lif !A)}$
  છે તે બતાવો.
\end{prob}

\end{document}
